\section{Thống kê mô tả}

\subsection{Thống kê mô tả cho các biến định lượng}
    Đầu tiên, ta tiến hành tính các đại lượng thống kê mô tả cơ bản (Mean, Median, SD, Min, Max) cho các biến định lượng, với trọng tâm phân tích là biến chi phí giao hàng (delivery\_charges).
    Ta sử dụng các tập lệnh (command):
    \begin{enumerate}
        \item function(x): tự định nghĩa hàm để tính 5 thông số thống kê.
        \item apply(X, MARGIN, FUN): áp dụng hàm tính toán vừa tạo lên toàn bộ các cột dữ liệu định lượng.
    \end{enumerate}
    Phần code lập trình xử lý các thông số thống kê:

    \begin{lstlisting}
    clean_data <- read.csv("clean_data_full.csv")

    describe_function <- function(x) {
      c(Mean = mean(x, na.rm = TRUE),
        Median = median(x, na.rm = TRUE),
        SD = sd(x, na.rm = TRUE),
        Min = min(x, na.rm = TRUE),
        Max = max(x, na.rm = TRUE))
    }
    num_vars <- c("delivery_charges", "coupon_discount", "order_total", "distance_to_nearest_warehouse")
    des_table <- apply(clean_data[, num_vars], 2, describe_function)
    print(des_table)
    \end{lstlisting}
    \textbf{Kết quả thu được:}\\
    Bảng thống kê mô tả cho thấy chi tiết giá trị trung bình, trung vị, độ lệch chuẩn, giá trị lớn nhất và nhỏ nhất của chi phí giao hàng cùng các biến liên quan.
    \begin{figure}[htbp]
        \centering
        \fbox{\includegraphics[width=0.9\linewidth]{hinh_1.jpg} }
        \caption{Kết quả bảng thống kê mô tả các biến định lượng.}
        \label{Bang_thong_ke_mo_ta}
    \end{figure}\\
    Mặc dù trong clean\_data thì customer\_lat và customer\_long là hai biến định lượng, nhưng ta không đưa vào bảng để tính 5 thông số thống kê cơ bản (Mean, Median, SD, Min, Max) vì:\\
    
        \textbf{Ý nghĩa thống kê thực tiễn:} Việc tính trung bình (Mean) hay độ lệch chuẩn (SD) của vĩ độ và kinh độ khách hàng không mang lại bất kỳ ý nghĩa phân tích kinh tế hay logistics nào.\\
        
        \textbf{Vai trò riêng:} Các biến customer\_lat và customer\_long không dùng để tính toán các đại lượng đặc trưng trung tâm hay phân tán, mà chỉ được giữ lại chủ yếu để kiểm tra sự phân bố không gian và vẽ các biểu đồ (thường liên quan đến biến distance\_to\_nearest\_warehouse).  
\subsection{Thống kê số lượng cho các biến phân loại}
    Tiếp theo, ta lập bảng thống kê số lượng (tần số) để xem xét sự phân bố đơn hàng theo các biến định tính: mùa vụ (season), trạng thái giao hàng nhanh (is\_expedited\_delivery), kho hàng (nearest\_warehouse) và mức độ hài lòng của khách hàng (is\_happy\_customer).
    Ta sử dụng lệnh:
    \begin{enumerate}
        \item table(): lập bảng tần số đếm số lượng phần tử thuộc mỗi nhóm phân loại.
    \end{enumerate}

    \begin{lstlisting}
    freq_season <- table(clean_data$season)
    freq_expedited <- table(clean_data$is_expedited_delivery)
    freq_warehouse <- table(clean_data$nearest_warehouse)
    freq_happy <- table(clean_data$is_happy_customer)
    
    print(freq_season)
    print(freq_expedited)
    print(freq_warehouse)
    print(freq_happy)
    \end{lstlisting}
    \textbf{Kết quả thu được:}\\
    Dữ liệu số lượng được phân bổ cụ thể cho từng hạng mục:\\
    
    \textbf{Mùa:}
    \begin{figure}[H]
        \centering
        \fbox{\includegraphics[width=0.9\linewidth]{hinh_2.png}}
        \caption{Kết quả số lượng theo mùa.}
        \label{Bang_thong_ke_mo_ta}
    \end{figure}
    
    \textbf{Nhận xét:} Dữ liệu cho thấy sự phân bố đơn hàng giữa các mùa có sự khác biệt. Cao nhất là mùa Xuân với 134 đơn hàng, tiếp đó là mùa Thu với 127 đơn, mùa Hè với 124 đơn và thấp nhất là mùa Đông với 115 đơn hàng. Từ đó cho thấy mức độ tiêu thụ ở mùa Xuân tăng cao sớm, trong khi đó lại giảm đi ở mùa Đông.\\
    
    \textbf{Trạng thái giao hàng nhanh:}
    \begin{figure}[H]
        \centering
        \fbox{\includegraphics[width=0.3\linewidth]{hinh_3.png}}
        \caption{Kết quả số lượng theo trạng thái giao hàng nhanh.}
        \label{Bang_thong_ke_mo_ta}
    \end{figure}
    
    \textbf{Nhận xét:} Sự phân bố đơn hàng theo trạng thái giao hàng nhanh cho thấy có 249 đơn yêu cầu giao hàng nhanh (TRUE) và còn lại 251 đơn không yêu cầu giao hàng nhanh (FALSE). Từ đó cho thấy xu hướng người mua không yêu cầu giao hàng nhanh cao, nhưng so với người mua yêu cầu giao hàng nhanh thì tỉ lệ này khá cân bằng, khẳng định vai trò của dịch vụ giao hàng nhanh.\\
    
    \textbf{Kho gần nhất:}
    \begin{figure}[H]
        \centering
        \fbox{\includegraphics[width=0.9\linewidth]{hinh_4.png}}
        \caption{Kết quả số lượng theo kho gần nhất.}
        \label{Bang_thong_ke_mo_ta}
    \end{figure}
    
    \textbf{Nhận xét:} Dữ liệu cho thấy sự phân bố đơn hàng giữa giữa các kho không đồng đều. Trong đó ta thấy Thompson chiếm số lượng nhiều nhất là 197 đơn, tiếp đó là Nickolson là 184 đơn và thấp nhất là Bakers với 119 đơn hàng. Điều này chỉ ra rằng kho Thompson phổ biến hơn cũng như chất lượng dịch vụ tốt hơn so với 2 kho còn lại.\\

    \textbf{Mức độ hài lòng:}
    \begin{figure}[H]
        \centering
        \fbox{\includegraphics[width=0.3\linewidth]{hinh_5.png}}
        \caption{Kết quả số lượng theo mức độ hài lòng.}
        \label{Bang_thong_ke_mo_ta}
    \end{figure}
    
    \textbf{Nhận xét:} Dữ liệu cho thấy đa số các khách hàng đều có đánh giá tốt với 359 đơn hàng với mức độ hài lòng cao (TRUE), nhiều hơn 2,5 lần so với số đơn có mức độ hài lòng thấp (FALSE).\\
\subsection{Mô tả dữ liệu bằng đồ thị}
    Sau khi ta thực hiện phân tích, tổng hợp về số lượng đơn hàng theo các biến nearest\_warehouse, season, is\_expedited\_delivery và is\_happy\_customer, để trực quan hóa sự phân bố và các mối liên hệ liên quan đến chi phí giao hàng (delivery\_charges), ta tiến hành vẽ các đồ thị chuyên sâu.

    \subsubsection{Biểu đồ phân phối tần suất (Histogram)}
        Ta sử dụng lệnh \texttt{hist()} để vẽ biểu đồ phân phối.
        \begin{lstlisting}
    hist(clean_data$delivery_charges, main="Histogram of Delivery Charges", col="lightblue", xlab="Delivery Charges (USD)")
        \end{lstlisting}

        \begin{figure}[H]
            \centering
            \fbox{\includegraphics[width=0.8\linewidth]{do_thi_1.png}}
            \caption{Biểu đồ Histogram của chi phí giao hàng.}
            \label{Bieu_do_Histogram_Barplot}
        \end{figure}
        \textbf{Nhận xét:} Biểu đồ Histogram cho thấy phần lớn chi phí giao hàng tập trung ở mức giá trung bình, phân phối có dạng hơi lệch phải. Từ đó ta kết luận được phần lớn là các đơn hàng có phí thấp, một số ít là các đơn hàng có phí cao. 

    \subsubsection{Biểu đồ cột theo kho hàng (Barplot)}
        Mặc dù đồ thị Barplot giúp đếm số lượng đơn hàng theo kho, nhưng nó là nền tảng để giải thích yếu tố nearest\_warehouse - một biến độc lập sẽ được đưa vào mô hình hồi quy để dự báo deliver\_charges ở phần sau.
        \begin{lstlisting}
    barplot(freq_warehouse, main="Orders by Warehouse", col="lightgreen")
        \end{lstlisting}

        \begin{figure}[H]
            \centering
            \fbox{\includegraphics[width=0.8\linewidth]{do_thi_2.png}}
            \caption{Biểu đồ Barplot của số lượng đơn hàng tại các kho.}
            \label{Bieu_do_Barplot}
        \end{figure}
        \textbf{Nhận xét:} Biểu đồ cho thấy số lượng đơn hàng tại cả ba kho. Lần lượt là kho Baker với 119 đơn, tiếp đến là kho Nickolson với 184 đơn và kho Thompson với 197 đơn, cao nhất trong 3 kho. Từ đồ thị, ta nhân thấy rỏ nhất rằng số lượng đơn tập trung chủ yếu ở kho Thompson và Nickolson, trong khi đó số đơn tại kho Bakers lại ít hơn nhiều. Từ đó ta dể dàng so sánh mức độ phân bố đơn hàng tại các kho. 

    \subsubsection{Biểu đồ hộp theo mùa (Boxplot)}
        Để kiểm tra sự phân phối của biến delivery\_charge theo season, ta sử dụng biểu đồ hộp (Boxplot) bằng lệnh boxplot().
        \begin{lstlisting}
    boxplot(delivery_charges ~ season, data=clean_data, main="Delivery Charges by Season", col=c("pink","lightblue","lightgreen","yellow"), ylab="Delivery Charges")
        \end{lstlisting}

        \begin{figure}[H]
            \centering
            \fbox{\includegraphics[width=1\linewidth]{do_thi_3.png}}
            \caption{Biểu đồ Boxplot theo mùa.}
            \label{Bieu_do_Tuong_quan}
        \end{figure}
        \textbf{Nhận xét:} Từ biểu đồ Boxpot, ta nhận thấy rằng chi phí vận chuyển vào mùa xuân và mùa hè cao hơn 2 mùa còn lại là mùa thu và mùa đông. Từ đó ta dễ dàng nhận thấy yếu tố mùa có ảnh hưởng đến chi phí vận chuyển.

    \subsubsection{Biểu đồ hộp theo trạng thái giao hàng (Boxplot)}
        Biểu đồ Boxplot theo trạng thái giao hàng trực tiếp đối chiếu mức phí giao hàng giữa nhóm khách yêu cầu giao hàng nhanh (TRUE) và không yêu cầu giao hàng nhanh (FALSE).
        \begin{lstlisting}
    boxplot(delivery_charges ~ is_expedited_delivery, data=clean_data, main="Delivery Charges by Expedited Status", col=c("pink","lightblue"), ylab="Delivery Charges")
        \end{lstlisting}

        \begin{figure}[H]
            \centering
            \fbox{\includegraphics[width=0.8\linewidth]{do_thi_4.png}}
            \caption{Biểu đồ Boxplot theo trạng thái giao hàng.}
            \label{Bieu_do_QQ_Plot}
        \end{figure}
        \textbf{Nhận xét:} Dựa vào biểu đồ Boxpot theo trạng thái giao hàng, ta nhận thấy chi phí giao hàng khi khách yêu cầu giao hàng nhanh cao hơn nhiều so với chi phí giao hàng khi khách không yêu cầu giao hàng nhanh.

        \subsubsection{Đồ thị ma trận tương quan (Correlation Matrix)}
        Ma trận tương quan giúp tính toán mối liên hệ tuyến tính trực tiếp giữa delivery\_charges với khoảng cách (distance\_to\_nearest\_warehouse) và vị trí (ccustomer\_lat, customer\_long), nhằm loại trừ hiện tượng đa cộng tuyến trước khi xây dựng mô hình dự báo. Ta sử dụng thư viện Corrplot. 
        \begin{lstlisting}
    library(corrplot)
    cor_vars <- c("delivery_charges", "customer_lat", "customer_long", "distance_to_nearest_warehouse")
    cor_matrix <- cor(clean_data[, cor_vars])
    corrplot(cor_matrix, method="number")
        \end{lstlisting}

        \begin{figure}[H]
            \centering
            \fbox{\includegraphics[width=0.8\linewidth]{do_thi_5.png}}            \caption{Đồ thị ma trận tương quan (Correlation Matrix)}
            \label{Bieu_do_QQ_Plot}
        \end{figure}
        \textbf{Nhận xét:} Dựa vào đồ thị ma trận tương quan, ta nhận thấy được hệ số tương quan ứng với từng cặp biến độc lập gần bằng 0. Từ đó ta suy ra được các biến độc lập không có quan hệ tuyến tính với nhau, dẫn đến không xảy ra hiện tượng đa cộng tuyến.

        \subsubsection{Biểu đồ phân tán (Scatter Plot)}
        Biểu đồ phân tán thể hiện mối quan hệ tuyến tính giữa delivery\_charges với các biến distance\_to\_nearest\_warehouse, customer\_lat và customer\_long. Nó làm cơ sở cho mô hình hồi quy tuyến tính. 
        \begin{lstlisting}
    plot(clean_data$distance_to_nearest_warehouse, clean_data$delivery_charges,
     main = "Scatter Plot: Distance vs Delivery Charges",
     xlab = "Distance to Nearest Warehouse",
     ylab = "Delivery Charges",
     col = "darkred", pch = 16)
    plot(clean_data$customer_lat, clean_data$delivery_charges,
     main = "Scatter Plot: Customer lat vs Delivery Charges",
     xlab = "Customer lat ",
     ylab = "Delivery Charges ",
     col = "darkred", pch = 16)
    plot(clean_data$customer_long, clean_data$delivery_charges,
     main = "Scatter Plot: customer long vs Delivery Charges",
     xlab = "Customer long ",
     ylab = "Delivery Charges ",
     col = "darkred", pch = 16)
        \end{lstlisting}

        \begin{figure}[H]
            \centering
            \fbox{\includegraphics[width=0.7\linewidth]{do_thi_6.1.png}}
            \caption{Biểu đồ phân tán giữa Delivery Charge và Distance to Nearest Warehouse}
            \label{Bieu_do_QQ_Plot}
        \end{figure}

        \begin{figure}[H]
            \centering
            \fbox{\includegraphics[width=0.7\linewidth]{do_thi_6.2.png}}
            \caption{Biểu đồ phân tán giữa Delivery Charge và Customer lat}
            \label{Bieu_do_QQ_Plot}
        \end{figure}

        \begin{figure}[H]
            \centering
            \fbox{\includegraphics[width=0.7\linewidth]{do_thi_6.3.png}}
            \caption{Biểu đồ phân tán giữa Delivery Charge và Customer long}
            \label{Bieu_do_QQ_Plot}
        \end{figure}
        \textbf{Nhận xét:} Dựa vào 3 biểu đồ phân tán giữa delivery\_charges với customer\_lat, customer\_long và distance\_to\_nearest\_warehouse, ta nhận thấy các điểm dữa liệu phân bố không đồng đều hoặc tập trung thành từng cụm, không có một xu hướng tuyến tính rõ ràng. Điều này phù hợp với hệ số tương quan tương ứng với từng cặp biến độc lập gần bằng 0 như ở đồ thị ma trận tương quan. Chứng minh rằng khoảng cách hoặc tọa độ thông thường không có quan hệ tuyến tính với chi phí giao hàng.

        \subsubsection{Biểu đồ kiểm tra phân phối chuẩn (QQ-Plot)}
        Biểu đồ QQ-Plot kiểm tra trực tiếp xem biến delivery\_charges có tuân theo phân phối chuẩn hay không.
        \begin{lstlisting}
    qqnorm(clean_data$delivery_charges, 
       main = "Normal Q-Q Plot of Delivery Charges", 
       col = "black")
       qqline(clean_data$delivery_charges, col = "red", lwd = 2)
        \end{lstlisting}

        \begin{figure}[H]
            \centering
            \fbox{\includegraphics[width=0.8\linewidth]{do_thi_7.png}}
            \caption{Biểu đồ QQ-Plot kiểm tra phân phối chuẩn.}
            \label{Bieu_do_QQ_Plot}
        \end{figure}
        \textbf{Nhận xét:} Dựa vào đồ thị Normal Q-Q Plot, ta quan sát thấy các điểm dữ liệu phần lớn bám sát đường thẳng kỳ vọng chuẩn màu đỏ ở khu vực trung tâm. Tuy nhiên, ở hai khu vực phần đuôi (đặc biệt là đuôi phải), các quan trắc nằm lệch rõ rệt khỏi đường chuẩn. Điều này cho thấy phân phối của biến chi phí giao hàng có hiện tượng lệch và nhiều khả năng không tuân theo phân phối chuẩn một cách hoàn hảo.